% =====================================================================
% Project 1: Numerical Solution of ODE Initial Value Problems
% Topic 3 — Chemical Kinetics: The Robertson Problem
% Team of 5 | Numerical Analysis of ODEs and PDEs
% =====================================================================
% HOW TO USE THIS TEMPLATE
% - Every section below maps 1:1 to a requirement in the Project Brief
%   (Background -> Model -> Methods -> Stability -> Numerical
%   experiments -> Conclusions) and to the topic-specific tasks in
%   Problem Pack Section 3.
% - Look for "% TODO:" comments: they tell you exactly what evidence
%   the grader expects in that spot. Delete the comment once filled.
% - Target length: >= 8 pages (excluding appendices/references).
% =====================================================================

\documentclass[11pt]{article}

% ---------------------- Packages ----------------------
\usepackage[utf8]{inputenc}
\usepackage[T1]{fontenc}
\usepackage{amsmath,amssymb,amsthm}
\usepackage{graphicx}
\usepackage{booktabs}
\usepackage{siunitx}
\usepackage{caption}
\usepackage{subcaption}
\usepackage{algorithm}
\usepackage{algpseudocode}
\usepackage{listings}
\usepackage{xcolor}
\usepackage[margin=2.5cm]{geometry}
\usepackage[colorlinks=true,linkcolor=blue,citecolor=blue,urlcolor=blue]{hyperref}
\usepackage{cleveref}

\lstset{
  basicstyle=\ttfamily\footnotesize,
  frame=single,
  breaklines=true,
  columns=fullflexible,
  language=Python
}

% ---------------------- Title ----------------------
\title{Numerical Solution of the Robertson Stiff Chemical Kinetics Problem:\\
Three Integration Methods, Stiffness Analysis, and Validated Long-Time Behaviour}
\author{
  Team \#XX \\
  Member 1 \and Member 2 \and Member 3 \and Member 4 \and Member 5 \\
  \small Numerical Analysis of ODEs and PDEs
}
\date{\today}

\begin{document}
\maketitle

\begin{abstract}
% TODO: 150-250 words. State: the phenomenon (autocatalytic chemical
% kinetics), why it is a canonical stiffness benchmark, the three
% methods compared, the two independent validation strategies used,
% and the headline quantitative result (e.g., the stiffness ratio
% growth and the explicit-Euler step restriction you *measured*,
% not just the literature estimate).
\end{abstract}

\tableofcontents
\newpage

% =====================================================================
\section{Background}
% =====================================================================
% TODO (Brief: "Background"):
% - Physical/chemical motivation: what real process does the Robertson
%   system schematically represent (autocatalytic reaction network),
%   why solvers for combustion/atmospheric chemistry care about this
%   class of problem.
% - Cite H.H. Robertson (1966) as the origin of the benchmark.
% - State the project question in "consultant" framing: given this
%   reacting system, which integrator can be trusted to resolve the
%   fast transient AND the slow relaxation over 8 orders of magnitude
%   in time, and how would you know if the answer is wrong?
% - One paragraph roadmap of the report.

% =====================================================================
\section{Model}
% =====================================================================

\subsection{Governing Equations}
% TODO: reproduce the ODE system and initial condition exactly as
% given in the Problem Pack -- do NOT introduce an extra factor of 2
% in the B+B term.
\begin{align}
  y_1' &= -0.04\, y_1 + 10^4\, y_2 y_3, \\
  y_2' &= 0.04\, y_1 - 10^4\, y_2 y_3 - 3\times 10^7\, y_2^2, \\
  y_3' &= 3\times 10^7\, y_2^2,
\end{align}
with $y(0) = (1,0,0)^T$, integrated on $t \in [0,40]$.

\subsection{Structural Properties}
% TODO:
% - Derive (1,1,1) f(y) = 0 explicitly (sum the three RHS terms) and
%   state the resulting invariant plane y1+y2+y3=1.
% - State the consequence: the full 3x3 Jacobian has a structural zero
%   eigenvalue for all y on the invariant plane. Show the eigenvector
%   direction if useful.
% - Introduce the reduced 2x2 system (eliminate y3 = 1-y1-y2) that you
%   will use for the stiffness-ratio analysis in Section 4.

\subsection{Reproducible Protocol}
% TODO: state verbatim the protocol you committed to (this must match
% what you actually did in Section 5):
% - Output grid: t=0 plus >= 200 log-spaced points in [1e-8, 40].
% - Reference solution: solver name (Radau/BDF/LSODA), tolerances,
%   and the "reduce tolerance by 100x, keep unchanged digits" check.
% - For uniform-step convergence studies: the >= 4 halved step sizes
%   used, and what "stable / non-negative / accurate" means operationally.

% =====================================================================
\section{Methods}
% =====================================================================

\subsection{Explicit Euler}
% TODO: one-step formula, order 1, local/global truncation error order.
\begin{equation}
  y_{n+1} = y_n + h\, f(t_n, y_n).
\end{equation}

\subsection{Classical Runge--Kutta (RK4)}
% TODO: Butcher tableau or stage formulas, order 4.

\subsection{Implicit Method: Implicit Euler}
% TODO: formula below; state why implicit Euler (L-stable, first
% order) was chosen over trapezoidal/CN for a stiff, strongly damped
% problem like Robertson (CN is only A-stable, not L-stable, and can
% ring on very stiff, non-oscillatory decay -- justify your choice).
\begin{equation}
  y_{n+1} = y_n + h\, f(t_{n+1}, y_{n+1}).
\end{equation}

\subsubsection{Newton's Method for the Nonlinear Stage Equation}
% TODO:
% - Write the residual F(y_{n+1}) = y_{n+1} - y_n - h f(t_{n+1}, y_{n+1}) = 0
% - Write the Newton update using the analytic Jacobian J = df/dy:
%   (I - h J(y^{(k)})) delta = -F(y^{(k)}),  y^{(k+1)} = y^{(k)} + delta
% - State your stopping criterion (residual norm < tol) and the
%   tolerance value used.
% - Report the tolerance-sweep experiment (Sanity check in Problem
%   Pack): show that Newton's tolerance is tight enough that algebraic
%   error is negligible relative to the time-discretisation error.
\begin{algorithm}[h]
\caption{Newton iteration for one implicit Euler step}
\begin{algorithmic}[1]
\State Initial guess $y^{(0)} \gets y_n$ (or predictor from previous step)
\For{$k = 0, 1, 2, \dots$}
  \State Evaluate $F(y^{(k)})$, $J(y^{(k)}) = \partial f/\partial y (t_{n+1}, y^{(k)})$
  \State Solve $(I - hJ(y^{(k)}))\,\delta = -F(y^{(k)})$
  \State $y^{(k+1)} \gets y^{(k)} + \delta$
  \If{$\|F(y^{(k+1)})\| < \text{tol}$} \textbf{break}
  \EndIf
\EndFor
\end{algorithmic}
\end{algorithm}

\subsection{Adaptive Step-Size Control}
% TODO: describe step-doubling OR an embedded pair; give the local
% error estimator formula and the step-acceptance/rejection rule
% (e.g. h_new = h * (tol/err)^(1/(p+1)) with safety factor).

\subsection{Implementation Notes}
% TODO (Grading Dimension C -- Implementation): explicitly state what
% you implemented yourselves vs. what you called from numpy/scipy
% (e.g., "we implemented Euler/RK4/implicit-Euler+Newton ourselves;
% we used scipy.integrate.solve_ivp(Radau) ONLY for the independent
% reference solution") and justify why. Confirm vectorisation.

% =====================================================================
\section{Stability and Stiffness Analysis}
% =====================================================================

\subsection{Absolute Stability Regions}
% TODO: derive and plot (single figure with 3 panels or 3 overlaid
% boundaries) the stability regions of Explicit Euler, RK4, and
% Implicit Euler in the complex h*lambda plane.
\begin{figure}[h]
  \centering
  % \includegraphics[width=0.8\textwidth]{figures/stability_regions.pdf}
  \caption{Absolute-stability regions of Explicit Euler, RK4, and Implicit Euler.}
  \label{fig:stability-regions}
\end{figure}

\subsection{Jacobian Eigenvalues and the Stiffness Ratio}
% TODO:
% - Plot both nonzero eigenvalues of the reduced Jacobian vs. log(t).
% - Plot the stiffness ratio S(t) = max|Re(lambda)| / min|Re(lambda)|
%   (nonzero eigenvalues only) vs. log(t), starting at the first
%   positive output time (NOT t=0).
% - State your numerical threshold for classifying an eigenvalue as
%   "zero".
% - Compare your computed values against the orientation values in the
%   Problem Pack (~3.0e3 at t=1e-4, ~5.4e3 at t=1e-2, ~1.58e5 at t=40)
%   -- do not simply quote them, show your own computed curve.
\begin{figure}[h]
  \centering
  % \includegraphics[width=0.8\textwidth]{figures/stiffness_ratio.pdf}
  \caption{Nonzero Jacobian eigenvalues and stiffness ratio $S(t)$ vs. $\log_{10} t$.}
  \label{fig:stiffness-ratio}
\end{figure}

\subsection{Practical Step-Size Restriction for Explicit Euler}
% TODO:
% - State the local frozen-Jacobian estimate h < 2/|lambda|_max and
%   its numeric value at t=40.
% - Explicitly flag this as a LOCAL diagnostic, not a nonlinear
%   stability proof (required statement per Brief Section 2).
% - Report the empirical step-size sweep experiment that tests this
%   prediction: run Explicit Euler at several h values and record
%   whether each run is stable / accurate / non-negative. Show
%   agreement or discrepancy with the theoretical estimate.

% =====================================================================
\section{Numerical Experiments}
% =====================================================================

\subsection{Reference Solution and Validation Chain}
% TODO: report all checks used per Brief Section 2:
% 1. Independent oracle (solve_ivp, Radau/BDF, rtol=1e-12) with the
%    tolerance-halving self-consistency check.
% 2. Known case / conserved quantity: the linear invariant
%    y1+y2+y3=1, tracked as a diagnostic (not a full validation on its
%    own -- see Section 5.3).
% 3. Self-consistency: agreement of RK4 and Implicit Euler in the
%    fine-mesh limit.
% State explicitly what tolerance/criterion you accepted and which
% digits of the reference you trust.

\subsection{Order of Convergence}
% TODO: error-vs-step-size log-log plot for at least two methods
% (recommend RK4 and Implicit Euler), using >= 4 halved step sizes,
% each run tagged stable/non-negative/accurate. Report the fitted
% slope vs. the theoretical order.
\begin{figure}[h]
  \centering
  % \includegraphics[width=0.7\textwidth]{figures/convergence_order.pdf}
  \caption{Error vs. step size (log-log) for RK4 and Implicit Euler, with fitted convergence order.}
  \label{fig:convergence}
\end{figure}

\subsection{Explicit Euler Failure Modes}
% TODO: three SEPARATE diagnostics, not conflated:
% (a) instability (blow-up / oscillation),
% (b) inaccuracy (large error even when bounded),
% (c) negativity (y_i < 0).
% Show the h at which each failure mode first appears, and compare to
% the theoretical bound from Section 4.3.

\subsection{Conservation, Positivity, and Accuracy: Three Separate Diagnostics}
% TODO: plot together but discuss separately:
% - |y1+y2+y3 - 1| vs t for each method.
% - min_i y_i(t) vs t (non-negativity).
% - full-state error ||y_numerical - y_reference|| vs t.
% Explicitly make the point: conservation alone does not imply
% accuracy (an inaccurate or even negative solution can still sum to 1).
\begin{figure}[h]
  \centering
  % \includegraphics[width=0.8\textwidth]{figures/three_diagnostics.pdf}
  \caption{Conservation defect, minimum component value, and full-state error, tracked separately.}
  \label{fig:diagnostics}
\end{figure}

\subsection{Implicit Euler vs. Explicit Euler at Matched Error}
% TODO: comparison table: step count / wall time / RHS+Jacobian
% evaluations needed by each method to reach a stated common error
% level. "Did not blow up" is explicitly NOT sufficient here.

\subsection{Adaptive Step-Size Control in Action}
% TODO: plot accepted step size vs. t (should shrink drastically near
% the fast transient and grow during slow relaxation), and plot
% estimated local error vs. the tolerance line.

\subsection{Final-Value Cross-Check}
% TODO: table of y2(t=40) from your implicit method at several
% tolerances, converging toward your own tightened reference; compare
% (do not just copy) against the Problem Pack orientation value
% y(40) ~ (0.7158271, 9.1855e-6, 0.2841637).
\begin{table}[h]
\centering
\begin{tabular}{@{}lccc@{}}
\toprule
Tolerance / $h$ & $y_1(40)$ & $y_2(40)$ & $y_3(40)$ \\
\midrule
% TODO: fill rows \\
\bottomrule
\end{tabular}
\caption{Convergence of the implicit-method solution at $t=40$ as tolerance/step size is refined.}
\label{tab:final-value}
\end{table}

% =====================================================================
\section{Conclusions}
% =====================================================================
% TODO:
% - Answer the project's mathematical question directly: which
%   method(s) can be trusted for this system, under what step-size
%   regime, and what is the quantitative evidence?
% - Summarise the stiffness ratio growth (your measured curve) and its
%   consequence for explicit methods.
% - State limitations: local frozen-Jacobian diagnostics vs. actual
%   nonlinear behaviour; what you did NOT verify.
% - One or two sentences on what a stronger (star3) extension would
%   look like (e.g., positivity-preserving high-order method), even if
%   not attempted.

% =====================================================================
\newpage
\appendix

\section{AI Transparency Log}
% TODO (required by Brief Section 8): for each significant AI-assisted
% step, record: what you asked for, what the AI produced, what you
% changed and why, and how you verified correctness (unit test,
% special case, cross-check against reference). Table format works
% well:
\begin{table}[h]
\centering
\small
\begin{tabular}{@{}p{2.5cm}p{4cm}p{4cm}p{3cm}@{}}
\toprule
Date/Step & AI Prompt (summary) & What we changed \& why & Verification \\
\midrule
% TODO: fill rows \\
\bottomrule
\end{tabular}
\caption{AI Transparency Log.}
\end{table}

\section{Individual Contribution Statement (ICS)}
% TODO: for each of the 5 team members, state concretely what they
% implemented, analysed, and wrote, and confirm they can explain the
% full project (this is cross-checked live in the Seminar Q\&A).

% =====================================================================
% \bibliographystyle{plain}
% \bibliography{references}
% TODO: at minimum cite H.H. Robertson (1966) and any solver/reference
% documentation used (e.g., scipy solve_ivp docs for Radau/BDF/LSODA).

\end{document}
